% GENERATED FILE. Edit canonical lessons, then run npm run generate:books.
% canonical-source-hash: fnv1a64:676c041c79e26652
% canonical-lessons: IT-C300-travolgere, IT-C300-ferito, IT-C300-ferire, IT-C300-illeso, IT-C300-cavarsela

\chapter{Feriti e illesi}
\label{ch:it-feriti-e-illesi}
\hlchaptermodality{300}

\begin{chapteropening}
\textbf{By the end of this chapter:} \emph{I can say who was hurt in an accident, who was run over and who came out unhurt.}

Complete the last lesson of chapter 300: I can say who was hurt in an accident, who was run over and who came out unhurt.
\end{chapteropening}

\section[travolgere]{travolgere — to run over, to sweep away}

\label{lesson:IT-C300-travolgere}



\begin{quote}
Before the new one: say the Italian for to dent, then the Italian for to damage.
\end{quote}

\subsection*{You'll want to know: travolgere}
\textbf{travolgere} — \textquotedblleft{}to run over, to sweep away\textquotedblright{}.

\begin{quote}
Un'auto ha travolto un ciclista.
\end{quote}

\begin{quote}
\emph{A car ran over a cyclist.}
\end{quote}

Past participle \emph{travolto}. Also a wave or a flood: \emph{travolto dall'acqua}.

\begin{cousinweb}
\emph{travolgere} joins \emph{tra-} \textquotedblleft{}across\textquotedblright{}, from Latin \emph{trāns}, and \emph{volgere} \textquotedblleft{}to turn\textquotedblright{}, from Latin \emph{volvere} \textquotedblleft{}to roll\textquotedblright{} (Lewis \& Short).
\end{cousinweb}

\begin{grammarlens}[title={What you've built}]
One more word for telling, explaining and describing at B1.
\end{grammarlens}

\subsection*{Your turn}
\begin{itemize}
\raggedright
  \item \emph{Say it:} \emph{travolgere}
  \item \emph{Say it:} \emph{travolgere}, once more
  \item \emph{Say it:} a sentence of your own with \emph{travolgere}
  \item \emph{Recall:} say the Italian for to dent, then the Italian for to damage, then say \emph{travolgere} again
\end{itemize}

\begin{tcolorbox}[breakable,colback=teal!4,colframe=teal!35!black,title={Before you move on}]
What does travolgere mean? (\textquotedblleft{}To run over, to sweep away\textquotedblright{}.) Say it once more.
\end{tcolorbox}
\section[il ferito]{il ferito — the injured person, the wounded}

\label{lesson:IT-C300-ferito}



\begin{quote}
Before the new one: say the Italian for to damage, then the Italian for to run over.
\end{quote}

\subsection*{You'll want to know: il ferito}
\textbf{il ferito} — \textquotedblleft{}the injured person, the wounded\textquotedblright{}.

\begin{quote}
I feriti sono in ospedale.
\end{quote}

\begin{quote}
\emph{The injured are in hospital.}
\end{quote}

Plural \emph{i feriti}; as an adjective, \emph{ferito} means injured.

\begin{cousinweb}
\emph{ferito} is the past participle of \emph{ferire} \textquotedblleft{}to wound\textquotedblright{}, which is inherited from Latin \emph{ferīre} \textquotedblleft{}to strike\textquotedblright{} (EtymDB).
\end{cousinweb}

\begin{grammarlens}[title={What you've built}]
One more word for telling, explaining and describing at B1.
\end{grammarlens}

\subsection*{Your turn}
\begin{itemize}
\raggedright
  \item \emph{Say it:} \emph{il ferito}
  \item \emph{Say it:} \emph{il ferito}, once more
  \item \emph{Say it:} a sentence of your own with \emph{il ferito}
  \item \emph{Recall:} say the Italian for to damage, then the Italian for to run over, then say \emph{il ferito} again
\end{itemize}

\begin{tcolorbox}[breakable,colback=teal!4,colframe=teal!35!black,title={Before you move on}]
What does il ferito mean? (\textquotedblleft{}The injured person, the wounded\textquotedblright{}.) Say it once more.
\end{tcolorbox}
\section[ferire]{ferire — to injure, to wound}

\label{lesson:IT-C300-ferire}



\begin{quote}
Before the new one: say the Italian for to run over, then the Italian for the injured person.
\end{quote}

\subsection*{You'll want to know: ferire}
\textbf{ferire} — \textquotedblleft{}to injure, to wound\textquotedblright{}.

\begin{quote}
Nessuno si è ferito, per fortuna.
\end{quote}

\begin{quote}
\emph{Nobody was hurt, luckily.}
\end{quote}

\emph{Ferirsi} — to hurt oneself. Present like \emph{capire}: \emph{ferisce}.

\begin{cousinweb}
\emph{ferire} is inherited from Latin \emph{ferīre} \textquotedblleft{}to strike\textquotedblright{} (EtymDB).
\end{cousinweb}

\begin{grammarlens}[title={What you've built}]
One more word for telling, explaining and describing at B1.
\end{grammarlens}

\subsection*{Your turn}
\begin{itemize}
\raggedright
  \item \emph{Say it:} \emph{ferire}
  \item \emph{Say it:} \emph{ferire}, once more
  \item \emph{Say it:} a sentence of your own with \emph{ferire}
  \item \emph{Recall:} say the Italian for to run over, then the Italian for the injured person, then say \emph{ferire} again
\end{itemize}

\begin{tcolorbox}[breakable,colback=teal!4,colframe=teal!35!black,title={Before you move on}]
What does ferire mean? (\textquotedblleft{}To injure, to wound\textquotedblright{}.) Say it once more.
\end{tcolorbox}
\section[illeso]{illeso — unhurt, unharmed}

\label{lesson:IT-C300-illeso}



\begin{quote}
Before the new one: say the Italian for the injured person, then the Italian for to injure.
\end{quote}

\subsection*{You'll want to know: illeso}
\textbf{illeso} — \textquotedblleft{}unhurt, unharmed\textquotedblright{}.

\begin{quote}
Il bambino è uscito illeso dall'incidente.
\end{quote}

\begin{quote}
\emph{The child came out of the accident unhurt.}
\end{quote}

\emph{Uscire illeso} — to come out unharmed.

\begin{cousinweb}
\emph{illeso} is a learned borrowing from Latin \emph{illaesus} \textquotedblleft{}unhurt\textquotedblright{}, from \emph{in-} \textquotedblleft{}not\textquotedblright{} and \emph{laesus}, the past participle of \emph{laedere} \textquotedblleft{}to hurt\textquotedblright{} (Lewis \& Short).
\end{cousinweb}

\begin{grammarlens}[title={What you've built}]
One more word for telling, explaining and describing at B1.
\end{grammarlens}

\subsection*{Your turn}
\begin{itemize}
\raggedright
  \item \emph{Say it:} \emph{illeso}
  \item \emph{Say it:} \emph{illeso}, once more
  \item \emph{Say it:} a sentence of your own with \emph{illeso}
  \item \emph{Recall:} say the Italian for the injured person, then the Italian for to injure, then say \emph{illeso} again
\end{itemize}

\begin{tcolorbox}[breakable,colback=teal!4,colframe=teal!35!black,title={Before you move on}]
What does illeso mean? (\textquotedblleft{}Unhurt, unharmed\textquotedblright{}.) Say it once more.
\end{tcolorbox}
\section[cavarsela]{cavarsela — to come through all right, to get by}

\label{lesson:IT-C300-cavarsela}



\begin{quote}
Before the new one: say the Italian for to injure, then the Italian for unhurt.
\end{quote}

\subsection*{You'll want to know: cavarsela}
\textbf{cavarsela} — \textquotedblleft{}to come through all right, to get by\textquotedblright{}.

\begin{quote}
È stato un brutto incidente, ma ce la siamo cavata.
\end{quote}

\begin{quote}
\emph{It was a bad accident, but we came through it.}
\end{quote}

An idiom: \emph{me la cavo}, \emph{se l'è cavata}. Also ‘to manage’: \emph{me la cavo con l'inglese}.

\begin{cousinweb}
\emph{cavarsela} is built on \emph{cavare} \textquotedblleft{}to dig out, to pull out\textquotedblright{}, which is inherited from Latin \emph{cavāre} \textquotedblleft{}to hollow out\textquotedblright{} (Lewis \& Short): to pull oneself out of trouble.
\end{cousinweb}

\begin{grammarlens}[title={What you've built}]
One more word for telling, explaining and describing at B1.
\end{grammarlens}

\subsection*{Your turn}
\begin{itemize}
\raggedright
  \item \emph{Say it:} \emph{cavarsela}
  \item \emph{Say it:} \emph{cavarsela}, once more
  \item \emph{Say it:} a sentence of your own with \emph{cavarsela}
  \item \emph{Recall:} say the Italian for to injure, then the Italian for unhurt, then say \emph{cavarsela} again
\end{itemize}

\begin{tcolorbox}[breakable,colback=teal!4,colframe=teal!35!black,title={Before you move on}]
What does cavarsela mean? (\textquotedblleft{}To come through all right, to get by\textquotedblright{}.) Say it once more.
\end{tcolorbox}
